\documentclass[instructions.tex]{subfiles}
\begin{document}
 
\chapter{Du détachement des richesses.}
 
IL ne faut pas désirer d’être riche, 1. parce que les richesses ne peuvent rendre l’homme content; 2. parce qu’elles lui sont même à charge; 3. parce que ses désirs sont inutiles et pernicieux.
 
\primo Les richesses ne peuvent rendre l’homme content. Si les pauvres sont à plaindre, les riches le sont encore plus. Les pauvres ne sont à plaindre que parce qu’ils sont sans patience; mais ils ne sont pas à plaindre parce qu’ils sont pauvres. On trouve bien plus de pauvres contens que de riches heureux. Eussiez-vous toutes les richesses de la terre, vous diriez comme Salomon : Tout n’est que vanité. Plus on a de bien, plus on est avide; plus on est avide, plus on est misérable\footnote{Divites eguerunt. Ps. 33. 11.}; La cupidité, toujours insatiable, n’est jamais satisfaite; plus on possède, plus on désire: comme l’hydropique, plus il boit, plus il est altéré. Quel contentement peut avoir dans ses richesses un homme qui ne sait pas les employer utilement? Une poignée de sable au fond de la mer me servira autant qu’une bourse d’or à laquelle je n’ose toucher. Un amas de grains et de meubles précieux dont je n’use pas, ne me sert pas plus que s’il était à un marchand des Indes.
 
Il n’est pas sans exemple de voir des gens qui n’osent se servir de leurs biens, et qui s’épargnent le nécessaire. Agir de la sorte, au préjudice des pauvres et de sa famille, c’est une conduite indigne de l’honnête homme et du chrétien, une sordide avarice qui rend un homme méprisable, qui le rend pauvre et misérable dans son abondance, et criminel en même temps.
 
Mais quand je me servirais de mes richesses, est-ce pour moi un avantage d’en avoir beaucoup? Je serai vêtu d’habits plus précieux; mais ces vêtemens me rendront-ils plus heureux? Je serai nourri plus délicatement; mais cette nourriture n’empêchera pas que je ne sois sujet à plus d’infirmités que ceux qui n’ont que du pain. Je serai honoré; mais ce sont mes richesses qu’on respectera, tandis qu’on se moquera de ma sotte vanité. On admire le plumage de certains oiseaux, tandis qu’on méprise l’oiseau qui le porte. J’aurai dans mes richesses la satisfaction de ne manquer de rien; mais pour combien de temps? Cette nuit peut-être on me demandera mon âme.
 
\secundo Les grandes richesses ne peuvent donc rendre l’homme content; elles lui sont même onéreuses et souvent nuisibles; elles ne traînent après elles qu’affliction d’esprit, dit le Sage. Elles fatiguent quand on les amasse; elles embarrassent quand on les possède; elles souillent quand on les aime; elles tourmentent quand on les perd\footnote{Possessa onerant, amata inquiinant, amissa cruciant. S. Bern.}.
 
Des épines entrelacées qu’on porte sur sa main, ne font aucun mal; mais, dès qu’on les serre avec force, leurs pointes meurtrières blessent cruellement. Les richesses, dit l’Evangile, sont des épines; dès qu’on s’y attache et qu’on les aime, que de blessures mortelles ne font-elles pas dans le cœur! l’homme de bien qui les méprise, jouit d’une paix que n’éprouvent point ceux qui en sont esclaves\footnote{ Melius est modicum justo, super divitias peccatorum multas. Ps. 36.}. Combien de gens ont avoué qu’ils avaient bien plus de repos étant pauvres qu’après être devenus riches! Pourvu que nous ayons la nourriture et le vêtement, dit saint Paul, soyons contens. De quoi sert le reste? Pourquoi nous empresser d’acquérir ce qui nous ôtera notre repos et notre tranquillité?
 
\tertio Chacun néanmoins désire les richesses, parce qu’on ne comprend pas le trésor renfermé dans la pauvreté et dans la médiocrité. Mais désirs inutiles! car, pourquoi voulez-vous être riche, puisque Dieu ne le veut pas et qu’il vous en ôte les moyens? Désirs pernicieux et funestes à l’âme! Ceux qui veulent devenir riches, dit saint Paul, tombent dans les filets du démon. Ce désir aveugle l’homme, il est la source de tous les maux\footnote{Radix enim omnium malorum est cupiditas. 1. Tim. 6.}. On veut être riche, et on veut l’être tôt; et le Saint-Esprit nous avertit que celui qui se hâte de s’enrichir n’est pas innocent. On veut être riche, et on veut l’être à quelque prix que ce soit. On emploie tout, friponnerie, usure, monopole, procès, violences, mensonges, usurpation, etc. On profite de tout. Les mauvaises années, les calamités qui portent la désolation dans le cœur des peuples, réveillent la cupidité, pour profiter de la nécessité des pauvres gens. Désir d’être riche, combien as-tu aveuglé de gens avides et fait de misérables!
 
On veut être riche, et on veut l’être beaucoup. L’homme intéressé n’est point content de ce qu’il a, et croit avoir besoin de tout ce qu’il n’a pas. Il ne peut voir sans jalousie et sans inquiétude les possessions d’un voisin; il voudrait tout engloutir. Ses biens, ses terres, son argent, sont comme une boue épaisse qu’il amasse autour de lui, dans laquelle il ensevelit son cœur\footnote{Aggravet contrà se densum lutum, Habac. 2. Humbert. Pensées}. Oh! qu’on est insensé d’attacher ainsi son cœur à la terre, et d’être esclave de sa propre cupidité!
 
A l’exemple de Salomon, ne demandez ni les richesses, ni la pauvreté. Si vous êtes pauvre, ne désirez pas d’être riche. Ces désirs, dit saint Paul, sont inutiles, nuisibles, et plongent l’homme dans la mort. Si vous êtes riche, modérez vos désirs : on est toujours malheureux quand on n’est pas content de ce qu’on possède.
 
\section{Résolutions}
 
\primo Persuadé que les richesses ne rendent pas heureux celui qui les possède, qu’on n’en jouit pas long-temps, puisqu’on ne les emporte pas en l’autre vie, et qu’elles peuvent devenir aisément une source féconde de péchés, et la cause de la damnation éternelle, je ne m’y attacherai point. 

\secundo J’en supporterai avec résignation la perte. 

\tertio Et je pardonnerai volontiers les torts qui me seront faits. 

O mon Dieu! ne permettez pas que mon cœur que vous avez formé pour vous, s’attache à d’autres objets au préjudice de l’amour que je vous dois.
 
\end{document}
